% =============================================================================
% CAPÍTULO 6 — SPRINT 5: Despliegue, Empaquetado y Cierre de Versión
% =============================================================================
\chapter{Sprint 5 — Despliegue, Empaquetado y Cierre de Versión}
\label{cap:sprint5}

El Sprint 5 cierra la versión 2.0.0 del producto: preparación de la base de datos para
producción, despliegue de los servicios e infraestructura, compilación optimizada del frontend y
publicación de las Release Notes integrales.\par

\textbf{Periodo de ejecución:} del lunes 25 de mayo al sábado 6 de junio de 2026 (2 semanas).
\textbf{Módulo entregado (Plan General):} \textit{Estadísticas, Adaptabilidad y Lanzamiento}.\par

\section{Plan Estratégico de Lanzamiento (Release Plan)}
\label{sec:s5_resumen}

El plan de lanzamiento coordina las cuatro capas hacia un único artefacto monolítico
desplegable. La tabla~\ref{tab:s5_alcance} resume las actividades por capa.

\begin{table}[!ht]
    \centering
    \renewcommand{\arraystretch}{1.35}
    \fontsize{9}{10.8}\selectfont\sffamily
    \begin{tabularx}{\textwidth}{|l|X|}
        \hline
        \rowcolor{colorPrimario}
        \textcolor{white}{\textbf{Capa}} & \textcolor{white}{\textbf{Actividad de lanzamiento}} \\ \hline
        Base de datos & Seeders de datos maestros y políticas de backups. \\ \hline
        \rowcolor{grisTecnico}
        Backend & Variables de entorno seguras y despliegue en contenedores. \\ \hline
        Frontend & Build optimizado, code splitting y distribución de assets. \\ \hline
        \rowcolor{grisTecnico}
        ISO 26515 & Release Notes y empaquetado de la base de conocimiento. \\ \hline
    \end{tabularx}
    \caption{Alcance del Sprint 5 por capa.}
    \label{tab:s5_alcance}
\end{table}

El módulo del Plan General para este cierre es \textbf{Estadísticas, Adaptabilidad y
Lanzamiento}: paneles de métricas para el profesional y el administrador, adaptación final a
móvil, cambio de idioma y modo claro/oscuro, y el despliegue productivo del artefacto monolítico.

\subsection{Objetivos del Incremento}
\label{ssec:s5_objetivos}
\begin{enumerate}[noitemsep]
    \item \textbf{Estadísticas:} dashboards de métricas y reportes para profesional y administrador.
    \item \textbf{Adaptabilidad:} i18n (es/en/pt), modo claro/oscuro y responsividad final.
    \item \textbf{Producción de datos:} \textit{seeders} maestros y políticas de \textit{backups}.
    \item \textbf{Despliegue:} empaquetado monolítico, contenedores y variables seguras.
    \item \textbf{Cierre de versión:} Release Notes 2.0.0 y empaquetado de la base de conocimiento.
\end{enumerate}

\subsection{Historias de Usuario Abordadas}
\label{ssec:s5_backlog}
\begingroup
\footnotesize
\renewcommand{\arraystretch}{1.35}
\begin{TablaBacklog}
    IDAR-002 & Dashboard de métricas (admin) & Med & 8 & 5 & Visión general de la plataforma. \\ \hline
    UX-002 & Cambio de idioma (i18n) & Med & 5 & 5 & Traducción sin recarga (es/en/pt). \\ \hline
    UX-005 & Modo claro/oscuro & Med & 3 & 5 & Cambio de tema persistente. \\ \hline
    REL-001 & Despliegue monolítico & Alt & 8 & 5 & JAR único FE+BE en producción. \\ \hline
    REL-002 & Backups y recuperación & Alt & 5 & 5 & Política automatizada de respaldos. \\ \hline
    REL-003 & Release Notes 2.0.0 & Med & 3 & 5 & Notas de versión publicadas. \\ \hline
\end{TablaBacklog}
\endgroup

\subsection{Ceremonias del Sprint}
\label{ssec:s5_ceremonias}
La tabla~\ref{tab:s5_ceremonias} resume las ceremonias del incremento de cierre
(25 may--6 jun 2026), incluida la \textit{Sprint Review} de fin de versión.

\begin{table}[!ht]
    \centering
    \renewcommand{\arraystretch}{1.3}
    \fontsize{9}{10.8}\selectfont\sffamily
    \begin{tabularx}{\textwidth}{|l|X|}
        \hline
        \rowcolor{colorPrimario}
        \textcolor{white}{\textbf{Ceremonia}} & \textcolor{white}{\textbf{Resultado}} \\ \hline
        Planning & Plan de lanzamiento, seeders, despliegue y métricas. \\ \hline
        \rowcolor{grisTecnico}
        Daily & Coordinación del empaquetado y la contenerización. \\ \hline
        Review (versión) & Demostración del despliegue 2.0.0 al cliente. \\ \hline
        \rowcolor{grisTecnico}
        Retrospective & Cierre del proyecto y lecciones aprendidas. \\ \hline
    \end{tabularx}
    \caption{Ceremonias Scrum del Sprint 5 (cierre de versión).}
    \label{tab:s5_ceremonias}
\end{table}

\clearpage
